\documentclass[11pt]{article}

\usepackage[margin=1.1in]{geometry}
\usepackage{amsmath,amssymb,amsthm}
\usepackage{stmaryrd}
\usepackage{mathtools}
\usepackage{booktabs}
\usepackage{array}
\usepackage{xcolor}
\usepackage[colorlinks=true,linkcolor=black,citecolor=black,urlcolor=blue]{hyperref}
\usepackage[expansion=false]{microtype}

\theoremstyle{plain}
\newtheorem{theorem}{Theorem}
\newtheorem{proposition}[theorem]{Proposition}
\newtheorem{lemma}[theorem]{Lemma}
\newtheorem{corollary}[theorem]{Corollary}
\newtheorem{conjecture}[theorem]{Conjecture}
\theoremstyle{definition}
\newtheorem{definition}[theorem]{Definition}
\newtheorem{remark}[theorem]{Remark}
\newtheorem{example}[theorem]{Example}

% --- typography: the rho calculus is NEVER written with the Greek letter ---
\newcommand{\rhoc}{\textsf{rho}}
\newcommand{\GSLT}{\mathbf{GSLT}}
\newcommand{\Kc}{\mathsf{K}}
\newcommand{\Gg}{\mathbb{G}}
\newcommand{\Lo}{\mathcal{L}}
\newcommand{\op}{\mathrm{open}}
\newcommand{\bd}{\mathrm{build}}
\newcommand{\bis}{\sim}
\newcommand{\Terms}{\mathcal{T}}

\title{\bfseries Every Constructor an Interaction\\[.35em]
\large Spatial structure made visible to bisimulation\\
in interactive graph-structured lambda theories}
\author{L.\ G.\ Meredith\\ \small F1R3FLY.io}
\date{August 7, 2026}

\begin{document}
\maketitle

\begin{abstract}
Spatial logics separate more than bisimulation does. This is not a defect and it is not news: Caires
observed two decades ago that the equivalence induced by a spatial logic for the $\pi$-calculus sits
strictly between structural congruence and strong bisimilarity, and the intensionality of such logics
was settled shortly afterwards by Sangiorgi and by Hirschkoff, Lozes and Sangiorgi. Yet the same
observation, restated, is routinely taken to refute any claim that a generated spatial--behavioural
logic is \emph{adequate} --- that its formulae separate exactly what interaction separates. The two
claims look incompatible and are not.

They are not incompatible because bisimilarity is not one relation. It is indexed by the theory that
supplies the contexts, and a theory can be enlarged. We give a uniform, purely syntactic enlargement.
Given an interactive graph-structured lambda theory $\Gg$ with interaction constructor $\Kc$, adjoin
for each term constructor $C$ a nullary atom $C_\bullet$, overload $\Kc$ to accept an atom against an
argument list, and adjoin the two rewrite rules that build $C$ from its parts and open $C$ into its
parts. Nothing else. In the resulting theory $\Gg^+$ every term former is an interaction, and since
bisimulation observes interactions it observes every term former: bisimilarity on $\Gg^+$ collapses
onto the equational theory of $\Gg$, which is where the spatial connectives already were. The two
disputed propositions are then the two endpoints of one monotone family, indexed by how much of the
signature one chooses to open.

This is folklore in the spatial logics community. Our contribution is not the phenomenon but its
handling: earlier results reach extensionality \emph{semantically}, by exhibiting models --- typically
distributed calculi with failures --- whose observational equivalence happens to be characterised by
a spatial logic. We reach it \emph{syntactically}, by a transformation of presentations that is
mechanical in the signature and chooses nothing, and therefore \emph{generally}, for any interactive
GSLT rather than for a fixed calculus. We record what the construction costs. It is idempotent
because the adjoined constructors are nullary; it is calibrated, landing on the equational theory
rather than on raw syntax, because the equations are kept as equations; and it is not
object-capability safe, since unrestricted opening destroys encapsulation. We describe the shape of
the capability discipline that repairs this, observe that its factory is the unique constructor
admitting no destructor, and do not develop it here.
\end{abstract}

\tableofcontents

\section{Introduction}
\label{sec:intro}

\subsection{The phenomenon}

A spatial--behavioural logic for a process calculus carries two layers. The behavioural layer is
modal: formulae speak of the steps a term can take. The structural, or spatial, layer speaks of how
the term is put together --- that it is a parallel composition of two parts with such-and-such
properties, that it is a send on a name, that it is empty.

The layers do not have the same separating power, and the direction of the gap surprises people the
first time they meet it. Behavioural equivalence is coarse by design: it identifies terms that no
experiment can tell apart. Structural predicates are finer, because they can ask questions no
experiment can pose. The standard example is as old as CCS. In a calculus with prefix, sum and
parallel composition,
\[
a.\mathbf{0} \mid b.\mathbf{0}
\qquad\text{and}\qquad
a.b.\mathbf{0} + b.a.\mathbf{0}
\]
are strongly bisimilar --- the expansion law is exactly the statement that they are --- and they are
not structurally congruent. The first satisfies the separating conjunction
$\neg\mathbf{0} \mid \neg\mathbf{0}$: it splits into two nonempty parts. The second does not. A logic
with a separating conjunction therefore distinguishes two terms that bisimulation identifies.

\subsection{The apparent contradiction}

Now put that observation next to a second one. When the logic is not designed but \emph{generated}
from the presentation of the calculus --- one structural connective per term former, one modality per
rewrite rule and redex position --- one wants an adequacy theorem: logical equivalence coincides with
the calculus's observational equivalence. Adequacy is what makes a generated logic worth generating.
Too coarse and there are observable distinctions the logic cannot state; too fine and the logic
invites its user to chase differences no experiment could settle.

So one wants to assert both of:
\begin{align}
\text{logical equivalence} \;&=\; \text{context-labelled bisimulation}, \tag{A} \\
\text{the spatial layer} \;&\subsetneq\; \text{context-labelled bisimulation}. \tag{B}
\end{align}
Read carelessly these are inconsistent, and a careful reader will say so. We have seen (B) offered as
a refutation of (A) more than once, by readers who knew the literature well enough to recognise (B)
and not well enough to know what answers it.

\subsection{The resolution, in one sentence}

Bisimilarity is parameterised by the theory that supplies the contexts. (B) is a statement about
bisimilarity in $\Gg$. (A) is a statement about bisimilarity in a theory $\Gg^+$ obtained from $\Gg$
by a construction that makes term formation itself an interaction. Once the parameter is exhibited,
(A) and (B) are the two endpoints of a monotone family, and there is nothing to reconcile.

The construction is the subject of \S\ref{sec:construction}. Its slogan is:

\begin{quote}
\itshape Make every term former an interaction. Bisimulation observes interactions. Therefore
bisimulation observes every term former.
\end{quote}

\subsection{What is new here, and what is not}
\label{sec:contribution}

The phenomenon is known. Section~\ref{sec:lineages} sets out three lineages of prior work, of which
the third --- extensionality of spatial observation --- already establishes that spatial logics need
not be counted intensional. We claim no priority over any of it.

What we add is a change of method, and the change buys generality. The prior extensionality results
are \emph{semantic}: they exhibit a calculus, typically a distributed one with local and remote
communication and partial failures, and show that its observational equivalence is characterised by a
spatial logic. The extra separating power comes from features of that calculus --- failures, in
particular --- which are modelling choices about a particular domain.

Our construction is \emph{syntactic}. It is a transformation of presentations. It reads the signature,
adjoins one nullary atom per constructor and two rules per constructor, and stops. It chooses nothing,
it is defined for any interactive GSLT rather than for one calculus, and because it is mechanical in
the presentation it composes with the machinery that generates the logic from the same presentation.
That last point is the one that matters for our purposes: if both the logic and the enlargement are
functions of the signature, then adequacy is the statement that two functions of the same argument
agree, rather than a coincidence to be checked calculus by calculus.

We also record three things about the construction that we have not seen written down: that it is
idempotent, and why (\S\ref{sec:idempotent}); that its calibration --- landing on the equational
theory rather than below it --- depends on a presentational choice that is easy to get wrong
(\S\ref{sec:calibration}); and that it is not object-capability safe (\S\ref{sec:security}).

\subsection{What this note does not do}

It does not develop the capability discipline of \S\ref{sec:security}. It does not treat weak
equivalences, quantitative refinements, or the interaction with cost accounting beyond the sketch of
\S\ref{sec:cost}. And it makes no claim that opening a term is a sensible thing to permit in a
deployed system. The enlarged theory is a measuring instrument. It exists to say what bisimulation
\emph{can in principle} see; it is not a habitat.

\section{Three lineages}
\label{sec:lineages}

The literature that bears on this contains three distinguishable strands, and much of the confusion
we are addressing comes from reading them as one. We set them out in the order they arose.

\paragraph{The observation.} Caires, studying an expressive spatial logic for the synchronous
$\pi$-calculus with recursion, gives coinductive and equational characterisations of the equivalence
the logic induces and locates it: strictly between structural congruence and strong bisimilarity
\cite{CairesObs}. The companion development with Cardelli \cite{CairesCardelliI,CairesCardelliII} is
explicit that formulae are taken to denote processes up to structural equivalence rather than up to
bisimulation, precisely because a tensor operator that distinguishes bisimilar processes is difficult
to reconcile with the latter. This is (B), stated by the people who first had to confront it.

\paragraph{The intensionality results.} Sangiorgi \cite{SangiorgiExtInt} and then Hirschkoff, Lozes
and Sangiorgi \cite{HLS02,HLS03} settle where the induced equivalence sits: for logics of this kind
it is essentially structural congruence. Lozes adds the elimination results --- that in the static
ambient logic the adjuncts contribute no expressive power \cite{LozesAdjunct}, later strengthened by
Dawar, Gardner and Ghelli through model-comparison games \cite{DGG} --- and the picture is completed
by Caires and Lozes \cite{CairesLozes}, who show that a minimal dynamic spatial logic already encodes
action quantification and action modalities, so that the spatial and modal fragments have the same
separating power, coinciding with structural congruence modulo permutation of actions. Lozes' thesis
\cite{LozesThesis} is the systematic account.

This strand says the logic is \emph{intensional}: it sees syntax. It is often read, wrongly, as
saying that nothing can be done about this.

\paragraph{The extensionality results.} Something can be done about it. Caires and Vieira
\cite{CairesVieira} address the tension directly, showing that there are natural models in which
extensional observational equivalences \emph{are} characterised by spatial logics, composition and
void included --- and concluding that spatial observations need not always be counted intensional,
even when expressive enough to describe structure. Their development is carried out in a minimalist
distributed calculus with local synchronous communication, local computation, asynchronous remote
communication, and partial failures. The additional separating power available to an observer in that
setting is what closes the gap. A related model-theoretic treatment is given by Tuosto and Vieira
\cite{TuostoVieira}; Caires' tutorial survey \cite{CairesSurvey} is the best single entry point to the
area.

\begin{table}[t]
\centering
\begin{tabular}{@{}p{3.1cm}p{5.2cm}p{5.2cm}@{}}
\toprule
\textbf{strand} & \textbf{question settled} & \textbf{what it does not settle} \\
\midrule
the observation & where the induced equivalence sits relative to $\bis$ and $\equiv$ &
whether the gap is intrinsic \\
intensionality & that the induced equivalence is essentially $\equiv$; which connectives are
eliminable & whether an observer could be given the missing power \\
extensionality & that in suitable models observational equivalence is characterised spatially &
whether this can be obtained uniformly from a presentation \\
\bottomrule
\end{tabular}
\caption{The three strands, and the question each leaves open. The present note addresses the
right-hand cell of the third row.}
\label{tab:lineages}
\end{table}

Table~\ref{tab:lineages} is the summary we wish had been available to us. A reader who knows only the
second strand will conclude that spatial logics are irreducibly intensional and that any adequacy
claim about a generated logic is confused. That conclusion does not follow, and the third strand is
where the answer lives.

\section{Preliminaries}
\label{sec:prelim}

We recall only what is used. For graph-structured lambda theories in general we refer to
\cite{Omnibus}.

\begin{definition}[GSLT]
A \emph{graph-structured lambda theory} is a triple $\Gg = (\Sigma, E, R)$ where $\Sigma$ is the
signature of a lambda theory --- so constructors may bind, and terms may contain variables --- $E$ is
a set of equations over $\Sigma$-terms, and $R$ is a set of rewrite rules, the graph of rewrites being
adjoined to the structure rather than the structure being enriched over graphs. We write $\Terms(\Sigma)$
for the terms and $=_E$ for the least congruence containing $E$.
\end{definition}

\begin{definition}[Interactive]
$\Gg$ is \emph{interactive} when $\Sigma$ carries a distinguished binary constructor $\Kc$, the
\emph{interaction constructor} or \emph{cut}, such that the left-hand side of every rule in $R$ is of
the form $\Kc(l_1, l_2)$ up to $=_E$.
\end{definition}

Parallel composition plays this role for CCS, the $\pi$-calculus, the ambient calculus and the \rhoc\
calculus; application plays it for the lambda calculus.

\begin{definition}[Context-labelled transitions]
For $P, P' \in \Terms(\Sigma)$ write $P \xrightarrow{(\rho, D)} P'$ when there are a rule
$\rho = (l \to r) \in R$, a one-hole context $D[-]$ and a substitution $\sigma$ with
$P =_E D[l\sigma]$ and $P' =_E D[r\sigma]$. The label records both the rule and the position.
\emph{Context-labelled bisimulation} $\bis_\Gg$ is the largest symmetric relation closed under
matching of these labelled steps.
\end{definition}

\begin{remark}[Notation]
The scientist note \cite{ScientistNote} writes $K_j$ for what we here call $D$, and $K$ for the cut.
The coincidence of letters there is not accidental --- for the \rhoc\ calculus the cut and the redex
context are built from the same constructor --- but it is confusing in a general setting, so we
separate them: $\Kc$ is a constructor, $D$ is a context.
\end{remark}

\begin{definition}[The generated logic]
$\Lo(\Gg)$ has formulae
\[
\varphi \;::=\; \top \;\mid\; \neg\varphi \;\mid\; \varphi \wedge \psi
\;\mid\; C(\varphi_1,\dots,\varphi_n) \;\mid\; \langle \rho, D \rangle \varphi
\]
with one \emph{structural} connective $C(-,\dots,-)$ per $n$-ary constructor $C \in \Sigma$ and one
\emph{behavioural} modality per rule and redex position. Satisfaction for the structural connectives is
\[
P \models C(\varphi_1,\dots,\varphi_n)
\iff
P =_E C(t_1,\dots,t_n) \text{ with } t_i \models \varphi_i \text{ for each } i,
\]
and for the modalities $P \models \langle \rho, D\rangle \varphi$ iff $P \xrightarrow{(\rho,D)} P'$
with $P' \models \varphi$. We write $=_{\Lo(\Gg)}$ for logical equivalence.
\end{definition}

Because $\Kc$ is typically subject to associativity, commutativity and a unit in $E$, the connective
$\Kc(\varphi, \psi)$ is a separating conjunction: it holds of $P$ exactly when $P$ admits \emph{some}
decomposition into two parts satisfying $\varphi$ and $\psi$ respectively. The existential over
decompositions is the whole difficulty. It is not a step, so no modality expresses it, and that is
where (B) comes from.

\begin{remark}[Which connectives are admitted]
Not every generated logic is adequate for anything, and users of a generator have good reason to
restrict which connectives are emitted --- typically to obtain a fragment whose model checking
terminates \cite{Omnibus}. We take the full structural layer as the default and return to the
restricted case in \S\ref{sec:dial}, where it becomes a parameter rather than a caveat.
\end{remark}

\section{The construction}
\label{sec:construction}

Fix an interactive GSLT $\Gg = (\Sigma, E, R)$ with cut $\Kc$, and assume $\Gg$ is Turing complete.

\subsection{Lists, and why Turing completeness is enough for them}

$\Kc$ is binary; constructors have arbitrary arity. To present the arguments of an $n$-ary constructor
to $\Kc$ we need to tuple them, so we assume list formers $\mathsf{nil}$ and
$\mathsf{cons}$ with the evident equations, writing $[t_1,\dots,t_n]$ for the obvious term.

\begin{remark}[The lists are not an import]
\label{rem:lists}
Turing completeness supplies them: any of the standard encodings will do, and which one is chosen does
not affect anything below, since all that is used is that $[\,\cdot\,]$ is injective up to $=_E$ and
that its components are recoverable by the same construction we are about to define. This is worth
saying because it is the \emph{only} thing Turing completeness is used for in this note, and it is
worth remembering because in \S\ref{sec:security} we will need something that Turing completeness does
\emph{not} supply.
\end{remark}

\subsection{Atoms, overloading, and two rules per constructor}

\begin{definition}[The canonical extension]
\label{def:ext}
The \emph{canonical extension} $\Gg^+ = (\Sigma^+, E, R^+)$ of $\Gg$ is given by:
\begin{itemize}
\item $\Sigma^+ = \Sigma \uplus \{\, C_\bullet : C \in \Sigma \,\}$, where each $C_\bullet$ is
\emph{nullary};
\item $\Kc$ is overloaded so that $\Kc(C_\bullet, [t_1,\dots,t_n])$ is well formed whenever $C$ is
$n$-ary;
\item $E$ is unchanged;
\item $R^+ = R \cup \{\, \bd_C, \op_C : C \in \Sigma \,\}$ where, for $C$ of arity $n$,
\begin{align*}
\bd_C &: \;\; \Kc(C_\bullet, [x_1,\dots,x_n]) \;\longrightarrow\; C(x_1,\dots,x_n), \\
\op_C &: \;\; C(x_1,\dots,x_n) \;\longrightarrow\; \Kc(C_\bullet, [x_1,\dots,x_n]).
\end{align*}
\end{itemize}
We call the rules of $R$ \emph{proper} and those of $R^+ \setminus R$ \emph{administrative}.
\end{definition}

Three comments on the shape of this definition.

First, the overloading of $\Kc$ rather than the introduction of a fresh application former is the
point of the whole construction. $\Kc$ is not an arbitrary constructor; it is the one the transition
relation is indexed on, the one that defines what it means for two things to meet. Putting $C_\bullet$
into interaction position re-expresses construction as communication, and communication is precisely
what a context-labelled transition system is built to observe. Nothing is being added to the
observational apparatus. Syntax is being moved into the part of the theory the apparatus already
covers.

Second, $\bd_C$ alone would be a term \emph{builder}, and a builder reads nothing it did not build:
$C_\bullet$ placed beside an opaque $P$ fires only if $P$ is already a list of the right length. It is
$\op_C$ that makes structure legible, since it applies at any redex position and therefore to terms
nobody holds.

Third, $\op_C$ and $\bd_C$ are converse. Neither the theory nor any term loses anything by their
presence; what is gained is that the syntax tree of a term becomes a reachable part of its transition
behaviour.

\begin{remark}[The equation-free presentation]
\label{rem:pairs}
One may prefer a presentation with no equations at all --- for instance to hand the theory to a tool
that does only rewriting. Then orient each equation of $E$ in both directions, adding both to the
administrative rules. The two presentations agree up to administrative steps. We keep $E$ as equations
because doing so makes the calibration of \S\ref{sec:calibration} immediate rather than a proof
obligation, and because it localises the source of nondeterminism: with $E$ retained, the branching of
$\op_C$ comes from $E$-matching, which is exactly the existential over decompositions that the
separating conjunction quantifies. It is worth stressing that these must be \emph{pairs}. Orienting
the equations in one direction only would let bisimilarity fall below $=_E$ onto raw syntax, which
overshoots: the structural connectives respect $=_E$, so a relation finer than $=_E$ is finer than the
logic and adequacy fails from the other side.
\end{remark}

\begin{proposition}[$\Gg^+$ is an interactive GSLT]
\label{prop:still}
$\Gg^+$ is a GSLT, is interactive with the same cut $\Kc$, and the inclusion
$\Terms(\Sigma) \hookrightarrow \Terms(\Sigma^+)$ reflects proper transitions: for
$P, P' \in \Terms(\Sigma)$ and $\rho \in R$, $P \xrightarrow{(\rho,D)} P'$ in $\Gg$ iff the same holds
in $\Gg^+$.
\end{proposition}

\begin{proof}
Every rule of $R^+$ has a left-hand side of the form $\Kc(-,-)$ up to $=_E$: those of $R$ by hypothesis,
$\bd_C$ by construction, and $\op_C$ because $C(x_1,\dots,x_n) =_E \Kc(C(x_1,\dots,x_n), u)$ for the
unit $u$ of $\Kc$ when $\Kc$ has one. Where $\Kc$ has no unit, replace $\op_C$ by
$\Kc(\mathrm{Op}_\bullet, [C(x_1,\dots,x_n)]) \to \Kc(C_\bullet, [x_1,\dots,x_n])$ with
$\mathrm{Op}_\bullet$ a further nullary atom; nothing below changes. Reflection of proper transitions
holds because $R^+ \supseteq R$ and the added constructors are not mentioned in any rule of $R$, so no
$R$-redex is created or destroyed by their presence.
\end{proof}

\section{Basic properties}
\label{sec:basic}

\subsection{Calibration}
\label{sec:calibration}

\begin{lemma}[Calibration]
\label{lem:calib}
If $P =_E Q$ then $P \bis_{\Gg^+} Q$.
\end{lemma}

\begin{proof}
Transitions are defined up to $=_E$, so $=_E$-equal terms have identical labelled transition sets.
\end{proof}

Trivial as stated, and stated anyway, because it is the property that the alternative presentation of
Remark~\ref{rem:pairs} has to work for, and because it is the half of the calibration that fixes the
\emph{lower} bound: the enlargement must not push bisimilarity below the equational theory, or it
overshoots the logic.

\subsection{Exposure}

\begin{lemma}[Exposure]
\label{lem:exposure}
Let $P \in \Terms(\Sigma^+)$ and let $C$ be $n$-ary. Then
\[
P \xrightarrow{(\op_C,\, [-])} \Kc(C_\bullet, [t_1,\dots,t_n])
\quad\text{iff}\quad
P =_E C(t_1,\dots,t_n).
\]
Consequently the set of $\op$-transitions available to $P$ at the empty context enumerates exactly the
top-level $E$-decompositions of $P$, one transition per decomposition, with the atom in the label
recording which constructor.
\end{lemma}

\begin{proof}
Immediate from Definition~\ref{def:ext} and the definition of context-labelled transitions, which
matches redexes up to $=_E$.
\end{proof}

This is the lemma that does the work, and it is worth pausing on what it says in the case that
motivated the whole difficulty. Take $C = \Kc$ with $E$ containing associativity, commutativity and
the unit. Then $P$'s $\op_\Kc$-transitions at the empty context are in bijection with the ways of
writing $P$ as $Q \mid R$ --- including $P \mid \mathbf{0}$ and $\mathbf{0} \mid P$, and every
rebracketing and permutation of a parallel product. The existential over splits in
$P \models \varphi \mid \psi$ has become an existential over transitions. That is the entire content of
the construction, and everything else is bookkeeping.

\subsection{The reconstruction theorem}

\begin{theorem}[Reconstruction]
\label{thm:main}
For closed $P, Q \in \Terms(\Sigma)$,
\[
P \bis_{\Gg^+} Q \iff P =_E Q .
\]
\end{theorem}

\begin{proof}
($\Leftarrow$) is Lemma~\ref{lem:calib}.

($\Rightarrow$) By induction on the size of $P$, where size is the number of constructor occurrences
in some (hence any, since $E$ is size-respecting on the presentations we consider) $=_E$-representative
of least size.

Suppose $P \bis_{\Gg^+} Q$ and $P =_E C(t_1,\dots,t_n)$ for some $C \in \Sigma$. By
Lemma~\ref{lem:exposure}, $P \xrightarrow{(\op_C, [-])} \Kc(C_\bullet, [t_1,\dots,t_n])$. Since
$P \bis_{\Gg^+} Q$, the term $Q$ must match that step with the \emph{same} label, so $\op_C$ is the
rule and the context is empty; by Lemma~\ref{lem:exposure} again, $Q =_E C(s_1,\dots,s_n)$ for some
$s_i$, and
\[
\Kc(C_\bullet, [t_1,\dots,t_n]) \;\bis_{\Gg^+}\; \Kc(C_\bullet, [s_1,\dots,s_n]).
\]
The atom $C_\bullet$ appearing in the label is what makes the constructors agree: distinct
constructors give distinct atoms and hence distinct labels, so no rule other than $\op_C$ could have
matched.

It remains to descend into the arguments. The list $[t_1,\dots,t_n]$ is itself built from
$\mathsf{cons}$ and $\mathsf{nil}$, so Lemma~\ref{lem:exposure} applies again at the contexts
$\Kc(C_\bullet, -)$ and its iterates. Matching those steps pairwise yields
$t_i \bis_{\Gg^+} s_i$ for each $i$, because the labels record position and therefore align the
arguments in order. Each $t_i$ is smaller than $P$, so the induction hypothesis gives $t_i =_E s_i$,
whence $P =_E C(t_1,\dots,t_n) =_E C(s_1,\dots,s_n) =_E Q$.
\end{proof}

\begin{remark}[Where the lambda theory is used]
If $C$ binds --- an input prefix, an abstraction --- then opening it exposes an argument that is not
closed. This is the one place where it matters that a GSLT is built on a lambda \emph{theory} rather
than on a first-order signature: the argument of a binding constructor is an abstraction, which is a
term of the theory, and the induction goes through on abstractions as it does on closed terms. For a
first-order signature the construction is more elementary and the restriction to closed terms in
Theorem~\ref{thm:main} may be dropped.
\end{remark}

\begin{remark}[Infinite branching]
A term with replication has infinitely many $E$-decompositions and therefore infinitely many
$\op_\Kc$-transitions. Nothing in the proof requires image-finiteness; bisimulation is defined for
arbitrary branching and the induction is on term size, not on transition depth.
\end{remark}

\subsection{Idempotence}
\label{sec:idempotent}

A referee's next move, on seeing Theorem~\ref{thm:main}, is to run the construction again. $\Gg^+$ has
term formers that $\Gg$ did not; generate the logic over $\Gg^+$ and there will be structural
connectives for them; are those not again finer than $\bis_{\Gg^+}$, and is one not on a ladder with no
top?

\begin{proposition}[Idempotence]
\label{prop:idem}
$\bis_{\Gg^{++}} \;=\; \bis_{\Gg^+}$, and $=_{\Lo(\Gg^+)}$ adds no separating power over
$=_{\Lo(\Gg)}$ on $\Terms(\Sigma)$.
\end{proposition}

\begin{proof}
The constructors adjoined by Definition~\ref{def:ext} are nullary. A nullary constructor has no
arguments to conceal: its structural connective is the atomic proposition ``is $C_\bullet$'', and by
Lemma~\ref{lem:exposure} that proposition is already determined by the $\bd_C$-transitions available
at the term, hence already behavioural in $\Gg^+$. Adjoining atoms and rules for the atoms therefore
yields nothing not already present, and by Theorem~\ref{thm:main} both sides equal $=_E$.
\end{proof}

This is not a side remark. It is the reason $(-)^+$ is a construction rather than the first stage of a
colimit, and it is a consequence of a design decision that could easily have gone the other way: had
we introduced a family of $n$-ary destructors, one per constructor, rather than nullary atoms plus an
overloaded cut, the ladder would be real.

\section{Consequences}
\label{sec:consequences}

\subsection{Adequacy}

We now import the second strand of \S\ref{sec:lineages}.

\begin{proposition}[Intensionality; known]
\label{prop:known}
For the spatial logics considered in \cite{SangiorgiExtInt,HLS02,HLS03,CairesLozes}, logical
equivalence coincides with structural congruence, possibly modulo a permutation of actions. In the
present notation: $=_{\Lo(\Gg)}\; = \;=_E$, whenever the target logic specification carries enough
boolean structure to form characteristic formulae.
\end{proposition}

We state this as an import rather than prove it. The hypothesis is not decorative: a generated logic
whose specification omits negation, or omits conjunction, will not separate to $=_E$, and adequacy
will fail for reasons that have nothing to do with the present construction \cite{Omnibus}.

\begin{corollary}[Adequacy in the extension]
\label{cor:adequacy}
Under the hypothesis of Proposition~\ref{prop:known}, for closed $\Sigma$-terms
\[
=_{\Lo(\Gg)} \;=\; =_E \;=\; \bis_{\Gg^+} .
\]
That is: the logic generated from $\Gg$ is adequate for context-labelled bisimulation in $\Gg^+$.
\end{corollary}

\begin{corollary}[Strictness, preserved]
\label{cor:strict}
$\bis_{\Gg^+} \;\subsetneq\; \bis_{\Gg}$ in general, so the spatial layer does strictly refine
bisimulation in $\Gg$.
\end{corollary}

\begin{example}[The witness]
\label{ex:ccs}
In CCS, $P = a.\mathbf{0} \mid b.\mathbf{0}$ and $Q = a.b.\mathbf{0} + b.a.\mathbf{0}$ satisfy
$P \bis_\Gg Q$ by the expansion law. In $\Gg^+$ they are separated at the first step: $P$ has an
$\op_\Kc$-transition at the empty context to $\Kc(\Kc_\bullet, [a.\mathbf{0}, b.\mathbf{0}])$, and the
only $\op$-transition $Q$ has at the empty context is $\op_+$, carrying a different atom in its label.
Correspondingly $P \models \neg\mathbf{0} \mid \neg\mathbf{0}$ and $Q \not\models
\neg\mathbf{0} \mid \neg\mathbf{0}$.
\end{example}

Corollaries~\ref{cor:adequacy} and~\ref{cor:strict} are (A) and (B). They are both true, of different
relations, and the difference is a parameter that was suppressed rather than a disagreement.

\subsection{The dial}
\label{sec:dial}

Nothing forces the construction to be run on the whole signature.

\begin{definition}[Partial extension]
For $\mathcal{C} \subseteq \Sigma$ let $\Gg^{+\mathcal{C}}$ be Definition~\ref{def:ext} with atoms and
rules adjoined only for $C \in \mathcal{C}$.
\end{definition}

\begin{proposition}[Monotonicity]
\label{prop:mono}
If $\mathcal{C} \subseteq \mathcal{C}'$ then $\bis_{\Gg^{+\mathcal{C}'}} \subseteq
\bis_{\Gg^{+\mathcal{C}}}$, with $\bis_{\Gg^{+\varnothing}} = \bis_\Gg$ and
$\bis_{\Gg^{+\Sigma}} = {=_E}$.
\end{proposition}

\begin{proof}
Adding rules only adds transitions that must be matched, so the largest bisimulation can only shrink.
The endpoints are Proposition~\ref{prop:still} and Theorem~\ref{thm:main}.
\end{proof}

\begin{corollary}[Relativised adequacy]
Let $\Lo_\mathcal{C}(\Gg)$ be the fragment of the generated logic whose structural connectives are
those of $\mathcal{C}$. Then $=_{\Lo_\mathcal{C}(\Gg)}$ is adequate for $\bis_{\Gg^{+\mathcal{C}}}$
under the hypothesis of Proposition~\ref{prop:known} relativised to $\mathcal{C}$.
\end{corollary}

This is the form in which the result is most useful, and it puts the earlier caveat about restricted
connective sets in its proper place. A user who dials the generator down to a finitely checkable
fragment is not thereby abandoning adequacy; they are moving along the family of
Proposition~\ref{prop:mono}, and adequacy holds at every point of it, against the bisimulation of the
correspondingly partial extension. The logic and the enlargement are indexed by the same parameter,
and adequacy is the statement that they agree.

A second dial gives the same family without touching the signature: fix $\Gg^{+\Sigma}$ and vary
instead which administrative labels an observer is permitted to see. Observing none recovers
$\bis_\Gg$; observing all recovers $=_E$. Whether the two dials give the same family in general we
have not checked.

\subsection{Cost}
\label{sec:cost}

We record a direction rather than a result. In a cost-accounted GSLT \cite{CostMonad} every rewrite
carries a price, and the administrative rules are rewrites like any other. Two consequences suggest
themselves.

Evaluating a structural formula of depth $d$ against a term corresponds to a path of $d$ nested
$\op$-steps. So the price of a structural observation, which in \cite{ScientistNote} is stipulated as a
monotone schedule in formula depth, is on this reading a \emph{step count} --- monotone in depth
because depth is what it counts. And the separating conjunction is expensive for a reason that is now
visible rather than asserted: confirming $\varphi \mid \psi$ requires searching the
$E$-decompositions, which is a walk over administrative transitions rather than a single step.

Whether the induced schedule agrees with the stipulated one, and what the exchange rate between proper
and administrative steps ought to be, we leave open.

\section{The construction is not capability-safe}
\label{sec:security}

We state this plainly rather than hedging it, because a reader who knows object-capability discipline
will see it in the first line of Definition~\ref{def:ext} and will want to know whether we noticed.

Unrestricted $\op$ destroys encapsulation. The rule applies at every redex position, so any subterm of
any term, held by anybody or by nobody, may be opened; and opening yields the arguments as data. In a
calculus where an opaque reference is the unit of authority, this is fatal: open the reference and
harvest whatever authority its contents carry. Confinement fails, and with it every argument that
depends on it.

The repair is known in outline. Mike Stay's suggestion is that constructors and destructors should
each take a capability token, so that the right to build and the right to open are held rather than
ambient. There are two readings of such a token --- a \emph{permission}, licensing the rule, or a
\emph{brand}, so that $\bd_C$ stamps the term and $\op_C$ checks the stamp --- and the second is the
more useful, since it makes legibility a property of the term rather than a global property of the
world, which is what any application to distributed systems will want. Note that the token must gate
the \emph{pair}: the calibration of Lemma~\ref{lem:calib}, in the equation-free presentation of
Remark~\ref{rem:pairs}, requires that an observer able to take an administrative step is able to take
its converse, and an asymmetric gate breaks that.

Tokens require a factory: a constructor that manufactures them. Two observations, and then we stop.

\begin{remark}[The factory is an import]
Remark~\ref{rem:lists} used Turing completeness to obtain lists. It cannot be used again here.
Unforgeability is not a claim about what can be computed but about what a context can \emph{write
down}, and in a signature of free constructors every term is writable by anybody. So a genuine source
of privacy must be adjoined --- restriction, in the general case, giving a factory of the shape
$\mathsf{let}\ c = \mathsf{freshCap}()\ \mathsf{in}\ P$, which is essentially $(\nu c)P$. Reflective
calculi, in which a name is a quoted process, may be able to discharge the import internally; whether
quotation \emph{originates} privacy or merely propagates it from a seed is a question we do not settle
here.
\end{remark}

\begin{remark}[The factory is the unique constructor with no destructor]
\label{rem:nodestructor}
Whatever the factory is, it is the one constructor to which the construction of
Definition~\ref{def:ext} must \emph{not} be applied. Give it an $\op$ rule and a token can be opened;
if a token can be opened it can be rebuilt by $\bd$; if it can be rebuilt it can be forged; and if it
can be forged nothing is a capability and the gate collapses. So the capability-safe version of the
construction is not ``the canonical extension, with a caveat'' but ``the canonical extension, with
exactly one exception, and the exception is what makes the rest of it safe''. We think this is the
right way to state unforgeability in this setting: a capability is the output of the unique
destructor-free constructor.
\end{remark}

Beyond this the design has real content --- whether the token is conserved or purchased, whether it is
consumed by use or held, how it interacts with a cost-accounting resource signature --- and we do not
pursue it. It is worth being clear about the scope of the omission. Nothing in
\S\S\ref{sec:construction}--\ref{sec:consequences} depends on it, because $\Gg^+$ is used there as an
instrument for locating an equivalence, not as a calculus in which anything runs. The moment one wants
to \emph{execute} in $\Gg^+$, the capability discipline stops being optional.

\section{What this settles, and for whom}
\label{sec:settles}

We collect the consequences for the framework this note was written to serve.

\paragraph{Adequacy claims survive contact with intensionality results.} A generated spatial--behavioural
logic may be asserted adequate, provided the assertion names the relation: bisimulation in the
extension, not in the base. Assertions that omit the index are not wrong so much as ambiguous, and the
ambiguity is what invites the objection.

\paragraph{Structural observation is not a second channel.} An architecture that lists structural
inspection alongside interaction as a separate mode of access is, on this reading, double-counting.
Structural inspection \emph{is} interaction --- with administrative redexes --- and what distinguishes
the modes is not which channel they use but what they cost and what they destroy. That is a better
distinction, since it is the one that does work downstream.

\paragraph{Spatial and behavioural addresses coincide.} In $\Gg^+$ a decomposition site is a redex
position. Constructions that need to identify a structural locus with a modal label --- and there are
several in \cite{ScientistNote} --- are therefore well typed rather than analogical. This does not by
itself supply an enabling theorem: a syntactic locus that is a redex position in $\Gg^+$ need not
correspond to a \emph{proper} step that is reachable in $\Gg$, and claims about reachability still
require their own argument.

\paragraph{Morphisms.} The category $\GSLT$ has bisimulation-preserving maps as morphisms
\cite{Omnibus}, and an assignment that depends on spatial structure is not functorial over them; the
usual repair is to narrow the morphism class by hand to maps preserving the generated logic. Under
Theorem~\ref{thm:main} the narrowing is not a narrowing: logic-preserving maps of $\Gg$ are exactly
bisimulation-preserving maps of $\Gg^+$. The repair becomes functoriality on the image of $(-)^+$.

\paragraph{An in-principle ceiling.} Adequacy in the extension is a claim about what an ideal observer
could separate, not about what any particular agent can afford. Arguments that move from adequacy to a
conclusion about an actual learner need the affordability premise stated separately; the fragment an
agent occupies is fixed by its budget, and \S\ref{sec:dial} is where the two meet.

\section{Open problems}
\label{sec:open}

\begin{enumerate}
\item \textbf{Canonicity as an adjunction.} We have called $(-)^+$ canonical on the grounds that it is
mechanical in the presentation and chooses nothing. Is it a left adjoint --- free on $\Gg$ among
extensions making $\Lo(\Gg)$ behavioural, hence initial among them? An adjunction would upgrade
``evident construction'' to ``forced construction'', which matters wherever the absence of a design
choice is doing argumentative work.

\item \textbf{Do the two dials agree?} \S\ref{sec:dial} gives a family indexed by a subset of the
signature and a second indexed by a set of observable administrative labels. Are they the same family?

\item \textbf{The capability-indexed family.} With Stay's tokens the relation $\bis_{\Gg^+}$ becomes a
family indexed by the tokens an observer holds, interpolating between $\bis_\Gg$ and $=_E$ with the
index an object of the calculus rather than a parameter of the metatheory. This is the sharpest form
of the resolution and we have not developed it.

\item \textbf{No capability amplification.} The security claim wants a theorem: one cannot open one's
way to a token one was not given. The statement has a circular smell --- unforgeability rests on the
gate, the gate rests on tokens --- and needs a well-foundedness argument.

\item \textbf{Derived pricing.} Does the step count of \S\ref{sec:cost} agree with the stipulated
schedules used in cost-accounted models, and under what exchange rate?

\item \textbf{Weak equivalences.} Everything here is strong. Administrative steps are natural
candidates for being made silent, and the resulting weak theory is presumably where the connection back
to $\bis_\Gg$ is cleanest.

\item \textbf{Relation to the semantic route.} \cite{CairesVieira} obtains extensionality from partial
failures. Is there a translation carrying that model into a partial extension of the present kind, so
that failures and administrative openings turn out to be two presentations of one mechanism?
\end{enumerate}

\section*{Acknowledgements}

The construction was described to the author's satisfaction in conversation with Mike Stay, whose
capability refinement is the subject of \S\ref{sec:security} and whose objection is the reason that
section exists.

\begin{thebibliography}{99}

\bibitem{CairesObs} L.\ Caires. \emph{Behavioral and spatial observations in a logic for the
$\pi$-calculus.} In Proc.\ FOSSACS 2004, LNCS 2987, Springer, 2004.

\bibitem{CairesCardelliI} L.\ Caires and L.\ Cardelli. \emph{A spatial logic for concurrency (Part I).}
Information and Computation 186(2), 2003.

\bibitem{CairesCardelliII} L.\ Caires and L.\ Cardelli. \emph{A spatial logic for concurrency
(Part II).} Theoretical Computer Science 322(3):517--565, 2004.

\bibitem{SangiorgiExtInt} D.\ Sangiorgi. \emph{Extensionality and intensionality of the ambient logics.}
In Proc.\ POPL 2001, ACM Press, 2001.

\bibitem{HLS02} D.\ Hirschkoff, E.\ Lozes and D.\ Sangiorgi. \emph{Separability, expressiveness and
decidability in the ambient logic.} In Proc.\ LICS 2002, IEEE Computer Society, 2002.

\bibitem{HLS03} D.\ Hirschkoff, E.\ Lozes and D.\ Sangiorgi. \emph{Minimality results for the spatial
logics.} In Proc.\ FSTTCS 2003, LNCS 2914, Springer, 2003.

\bibitem{LozesAdjunct} E.\ Lozes. \emph{Adjunct elimination in the static ambient logic.} In Proc.\
EXPRESS 2003, ENTCS, Elsevier.

\bibitem{LozesThesis} E.\ Lozes. \emph{Expressivit\'e des logiques spatiales.} Th\`ese de doctorat,
ENS Lyon, 2004.

\bibitem{CairesLozes} L.\ Caires and E.\ Lozes. \emph{Elimination of quantifiers and undecidability in
spatial logics for concurrency.} In Proc.\ CONCUR 2004, LNCS 3170, Springer, 2004.

\bibitem{DGG} A.\ Dawar, P.\ Gardner and G.\ Ghelli. \emph{Adjunct elimination through games in static
ambient logic.} In Proc.\ FSTTCS 2004, LNCS 3328, Springer, 2004.

\bibitem{CairesVieira} L.\ Caires and H.\ T.\ Vieira. \emph{Extensionality of spatial observations in
distributed systems.} In Proc.\ EXPRESS 2006, ENTCS 175(3):131--149, Elsevier, 2007.

\bibitem{TuostoVieira} E.\ Tuosto and H.\ T.\ Vieira. \emph{An observational model for spatial logics.}
ENTCS 142:229--254, 2006.

\bibitem{CairesSurvey} L.\ Caires. \emph{Dynamic spatial logics: a tutorial survey.} Bulletin of the
EATCS 94, 2008.

\bibitem{Omnibus} L.\ G.\ Meredith. \emph{Graph-structured lambda theories: a reference for the working
developer.} F1R3FLY.io, 2026.

\bibitem{ScientistNote} L.\ G.\ Meredith. \emph{The mortal scientist: learning as cost-accounted
computation in a reflective higher-order calculus.} F1R3FLY.io, 2026.

\bibitem{ClassifierNote} L.\ G.\ Meredith. \emph{Classifying knot diagrams with a spatial--behavioural
logic: a logic obtained for free from the encoding of knots as processes.} F1R3FLY.io, 2026.

\bibitem{CostMonad} L.\ G.\ Meredith. \emph{Cost accounting as a monad for continued graph-structured
lambda theories.} F1R3FLY.io, 2026.

\end{thebibliography}

\end{document}
